
\begin{exercise}[subtitle={Proof of Hölder's inequality \Coffeecup}]
    \label{ex: hölder's inequality}
    Let \(\frac1p + \frac1q = 1\) with \(p,q > 1\).
    \begin{enumerate}[label=(\alph*)]
        \item\label{it: young inequality}
        Prove ``Young's inequality'': for \(a,b \ge 0\) we have
        \[
            ab \le \tfrac{a^p}p + \tfrac{b^q}q.
        \]
        \textbf{Hint:} \emph{Option 1}: Let \(x = p \log(a)\) and \(y = q \log(b)\)
        and use convexity of the exponential function.
        \emph{Option 2 (brute force)}: Find the minimum of 
        the function \(f(x) = \frac{x^p}p + \frac{b^q}q - xb\).
        \begin{solution}
        \end{solution}

        \item\label{it: holder's inequality proof}
        Prove ``Hölder's inequality'' 
        For \(p,q >1\) with \(\frac1p + \frac1q = 1\) we have
        \[
            \E{\abs{XY}} \le \E{\abs{X}^p}^{\frac1p}\E{\abs{Y}^q}^{\frac1q} = \norm{X}_{L^p} \norm{Y}_{L^q}.
        \]

        \textbf{Hint:} 
        Assume without loss of generality \(X \in \cL^p(\real)\) and \(Y\in \cL^q(\real)\),
        otherwise the right hand side is infinite and the claim is trivial.
        Apply Young's inequality to \(a =
        \frac{\abs{X}}{\norm{X}_{L^p}}\) and \(b = \frac{\abs{Y}}{\norm{Y}_{L^q}}\).
    \end{enumerate}
\end{exercise}

\begin{solution}
\begin{enumerate}[wide=0pt]
    \item[\ref{it: young inequality}]
    \emph{Option 1:} Using the convexity of the exponential function,
    \[\begin{aligned}
        ab
        = e^{\log(a) + \log(b)}
        = e^{\frac1p (p\log(a)) + \frac1q (q\log(b))}
        &\le \frac1p e^{p\log(a)} + \frac1q e^{q\log(b)}
        \\
        &= \frac{a^p}p + \frac{b^q}q.
    \end{aligned}
    \]
    \emph{Option 2:} On \([0, \infty)\), \(f\) is convex  as \(f''(x) = (p-1)x^{p-2} \ge 0\).
    The minimum is therefore attained at \(f'(x) = x^{p-1} - b = 0\), that is
    at \(x_* = b^{\frac1{p-1}}\). Since \(q= \frac{p}{p-1}\) we have that \(x_*^p = b^q\).
    \[
        f(x_*) = \frac{b^q}p + \frac{b^q}q - b^{\frac{1}{p-1} + 1} = 0.
    \]
    Consequently, for all \(x \ge 0\) we have \(f(x) \ge f(x_*) = 0\).
    In particular for \(x = a\) we have the claim.

    \item[\ref{it: holder's inequality proof}]
    Applying Young's inequality we have
    \[
        \frac{\abs{X}}{\norm{X}_{L^p}} \cdot \frac{\abs{Y}}{\norm{Y}_{L^q}}
        \le \frac1p \Bigl(\frac{\abs{X}}{\norm{X}_{L^p}}\Bigr)^p
        + \frac1q \Bigl(\frac{\abs{Y}}{\norm{Y}_{L^q}}\Bigr)^q.
    \]
    Taking expectation on both sides we get
    \[
        \frac{\E[\big]{\abs{XY}}}{\norm{X}_{L^p} \norm{Y}_{L^q}}
        \le 
            \frac1p \E[\Big]{\Bigl(\frac{\abs{X}}{\norm{X}_{L^p}}\Bigr)^p}
            + \frac1q \E[\Big]{\Bigl(\frac{\abs{Y}}{\norm{Y}_{L^q}}\Bigr)^q}
        = 
            \frac1p + \frac1q
        = 1,
    \]
    where we use that \(\norm{X}_{L^p}^p = \E{\abs{X}^p}\) is
    deterministic to move it outside of the expectation. Multiplying
    both sides by \(\norm{X}_{L^p} \norm{Y}_{L^q}\)
    implies the claim.
    \qedhere
\end{enumerate}
\end{solution}
